\begin{enumerate}[a.]
    \item Dos departamentos con el nombre 'Sistemas'  podrían existir al mismo tiempo.
    
    Se cumple. La única llave de \texttt{Departamento} es \texttt{NúmDepto}, \texttt{NombreDepto} no es clave ni único, por lo que puede repetirse. Pueden existir dos tuplas con el mismo \texttt{NombreDepto} pero diferente \texttt{NúmDepto}.

    \item Dos o más empleados pueden administrar el mismo \texttt{Departamento} al mismo tiempo.
    
    Se cumple. La llave primaria de \texttt{Administrar} está compuesta por \texttt{NúmEmpleado}, \texttt{NúmDepto} y \texttt{Desde}, por lo que dos empleados pueden administrar el mismo \texttt{Departamento} al mismo tiempo, pues sus \texttt{NúmEmpleado} serán diferentes y por lo tanto, las claves serán distintas.

    \item Un empleado puede trabajar en un \texttt{Departamento} y administrar otro al mismo tiempo.
    
    Se cumple. Las tablas \texttt{Trabajar} y \texttt{Administrar} son independientes y no existe ninguna restricción entre ellas, por lo que no se requiere que el departamento administrado sea el mismo en el que se trabaja.

    \item Para administrar un Departamento un empleado debe trabajar en dicho departamento.
    
    No se cumple. Aunque ambas tablas \texttt{Trabajar} y \texttt{Administrar} contienen llaves \texttt{NúmEmpleado} y \texttt{NúmDepto}, estas solo identifican al empleado y departamento, pero no requieren que cada tupla de \texttt{Administrar} tenga una correspondiente en \texttt{Trabajar}, por lo que el esquema permite administrar un departamento en el que no se trabaja.

    \item Un empleado podría trabajar en dos Departamentos a partir de la misma fecha.
    
    Se cumple. La llave primaria de \texttt{Trabajar} está compuesta por \texttt{(NúmEmpleado, NúmDepto, Desde)}, por lo que si un empleado trabaja en dos departamentos distintos, las tuplas tendrán diferentes \texttt{NúmDepto} y ambas serán válidas.

    \item Para las tuplas de la relación Administrar, hasta no puede ser anterior a desde.
    
    No se cumple. Aunque se define el periodo [desde,hasta), esto no impone una restricción que impida que \texttt{Hasta} sea anterior a \texttt{Desde}, además, se indica no considerar restricciones adicionales a las de cardinalidad y participación.

    \item Dado un empleado, podemos identificar exactamente el Departamento donde trabaja.
    
    No se cumple. La relación \texttt{Trabajar} tiene una llave primaria compuesta \texttt{(NúmEmpleado, NúmDepto, Desde)}, permitiendo que un mismo empleado tenga varias tuplas en \texttt{Trabajar}, además, por la participación parcial, podría incluso no tener ninguna, por lo que no podríamos identificar exactamente el \texttt{Departamento} donde trabaja.

    \item Ningún empleado puede cobrar más de un Salario al mismo tiempo.
    
    No se cumple. La tabla \texttt{Salario} tiene como llave primaria \texttt{(NúmEmpleado, Desde)}, lo que impide que un mismo empleado tenga dos tuplas con el mismo valor de \texttt{Desde}, sin embargo, no existe una restricción que impida salarios que se traslapen, por lo que sí podría cobrar más de un salario al mismo tiempo.

    \item Algunas tuplas en Salario podrían no tener valor para el atributo desde y ningún empleado asociado a ellas.
    
    No se cumple. La llave primaria de \texttt{Salario} es \texttt{(NúmEmpleado, Desde)}, por lo que ninguno de los dos valores puede ser nulo, además de esto, la participación del lado de \texttt{Salario} con \texttt{Empleado} es total.

    \item Un Departamento siempre tiene algún empleado que lo administre.
    
    No se cumple. El esquema permite que haya un \texttt{Departamento} sin algún empleado que lo administre, porque en la línea \texttt{Departamento} - \texttt{Administrar}, el lado de \texttt{Departamento} tiene participación parcial.

\end{enumerate}
